\section{Conclusions}

This paper introduced MouseBrainBench, a framework for evaluating the evidence
behind partial digital models of the mouse brain. The proposal was motivated by
a precise methodological problem: predictive correlation, reproducibility,
anatomical specificity, and mechanistic identifiability are frequently treated
as if they were interchangeable. MouseBrainBench represents them as separate
evidence blocks and makes the final mechanistic decision conjunctive.

The framework was evaluated under synthetic truth, a real negative control,
modern predictive benchmarks, and local functional connectomics. The synthetic
systems confirmed that reproducibility, topology, or direction alone cannot
pass the complete MIS. Allen Visual Behavior Neuropixels showed strong
reproducibility while failing topology-specific and directed-identifiability
requirements. Static and Dynamic Sensorium showed useful predictive signals but
did not provide the missing mechanistic evidence.

MICRONS supplied the main positive empirical result. Across 2991 co-registered
units in a discovery cohort and two non-overlapping hold-outs, connected pairs
showed greater readout-location similarity than distance- and degree-matched
non-connected pairs. Unit-cluster bootstrap intervals remained positive in all
cohorts. This reproduces a local observational structure-function association
compatible with stronger previously published MICRONS analyses; it is an
external-validity case for the audit, not a claim of biological priority.

The contribution is consequently narrower, and more defensible, than a digital
mouse brain. MouseBrainBench provides a reproducible mechanism for deciding
which interpretations survive the available controls. Future work should add
interventional datasets, test the framework in other circuits, formalize
claim-specific evidence blocks with external preregistration, and evaluate
external simulators under the same audit. A complete digital-twin claim should
remain outside scope until entity-specific, multiscale, and perturbational
evidence can be demonstrated together.
